\SourceLeaf{397}%
be the midpoints of \(AC,BD\) and \(EF\), respectively. Choosing the axis of abscissas and the origin entirely at will, I denote the coordinates of these nine points by \(a,a'\quad b,b'\quad c,c'\quad d,d'\) and so on. Since the three points \(A,B,E\) lie on a straight line, their coordinates satisfy the following condition
\[
a(e'-b')+b(a'-e')+e(b'-a')=0
\]
Likewise, since \(ADF,BCF,DCE\) are straight lines, we have
\[
\begin{aligned}
a(f'-d')+d(a'-f')+f(d'-a')&=0\\
b(c'-f')+c(f'-b')+f(b'-c')&=0\\
c(e'-d')+d(c'-e')+e(d'-e')&=0
\end{aligned}
\]

Adding these four equations gives
\[
(a+c)(e'+f'-b'-d')+(b+d)(a'+c'-e'-f')+(e+f)(b'+d'-a'-c')=0
\]
or, since clearly
\[
\begin{array}{lll}
\tfrac12(a+c)=g,&\tfrac12(b+d)=h,&\tfrac12(e+f)=i\\
\tfrac12(a'+c')=g',&\tfrac12(b'+d')=h',&\tfrac12(e'+f')=i'
\end{array}
\]
we get
\[
g(i'-h')+h(g'-i')+i(h'-g')=0
\]
This is the condition for \(G,H,I\) to lie on one straight line.
\par\bigskip\centerline{\rule{.25\textwidth}{.3pt}}
